\documentclass[conference]{IEEEtran}
\IEEEoverridecommandlockouts
% The preceding line is only needed to identify funding in the first footnote. If that is unneeded, please comment it out.
\usepackage{cite}
\usepackage{amsmath,amssymb,amsfonts}
\usepackage{algorithmic}
\usepackage{graphicx}
\usepackage{textcomp}
\usepackage{xcolor}
\def\BibTeX{{\rm B\kern-.05em{\sc i\kern-.025em b}\kern-.08em
    T\kern-.1667em\lower.7ex\hbox{E}\kern-.125emX}}
\begin{document}

\title{From PDF to Knowledge graph:\\ A Fine-tuned LLM vs RAG-KG vs Zero-shot LLM comparison \\
{\footnotesize Natural Language Processing and Text Mining - Project 24} \\
{\small \textbf{Code available at:} \texttt{https://github.com/raphsa/pdf\_to\_KnowledgeGraph\_LLM}}
}

\author{\IEEEauthorblockN{Raffaele Sali}
\IEEEauthorblockA{rsali25@student.oulu.fi}
\and
\IEEEauthorblockN{Martina Vianello}
\IEEEauthorblockA{mvianel25@student.oulu.fi}
\and
\IEEEauthorblockN{Beatrice Zani}
\IEEEauthorblockA{bzani25@student.oulu.fi}
}

\maketitle

\begin{abstract}
*CRITICAL: Do Not Use Symbols, Special Characters, Footnotes, 
or Math in Paper Title or Abstract.
\end{abstract}

\begin{IEEEkeywords}
component, formatting, style, styling, insert
\end{IEEEkeywords}

\section{Introduction}
Perform a concise short literature review (around one page) about the context and topic associated to your project with relevant literature, commented and used to introduce potential novelty of the topic of your project. It should also highlight the problem your project is attempting to solve and its potential impact. It is typical to view the various project specification as potential challenges that you try to solve in the methodology section. 

\section{Methodology}
In this section you provide the theoretical foundation and NLP pipeline that you will be using to solve the technical challenges.. Feel free to divide this section into subsection if needed. Draw high level (block diagram) and pseudo-code that describe the implementation details. Put a link to your Github account where dataset and code should be located. 
Avoid putting results in this section..

\section{Results and Discussions}
Put the graphical illustrations, tables and potential analysis or discussion of the obtained results in this section.
You can also divide this section into as many subsections as you want if you to report individual specification results in the corresponding subsection. Discussion should take into account the limitation of the data collection or provided data as well inherent limitation of the NLP pipeline or theoretical framework / hypothesis

\section{Overall Discussions}
You should elaborate on the overall findings of your project and potential link with state of art raised in the introduction section. Report also any subsequent work that might be needed to perform specification but cannot be done in the current task
Report any comparison results, if available
Report what you think you have performed as novel and worth pursuing.

\section{Conclusion}
Main conclusion of your work, difficulty of the tasks performed, skills gained and potential ways forward

\section*{References}
Follow homogeneous template of referencing

\section*{Appendix}
Report here any large graphical illustrations for instance or any extra documentation that you think is more suitable to be put in appendix

\end{document}
